\documentclass{article}
\usepackage{PRIMEarxiv} % Core package for the PrimeAI Template
\usepackage{amsmath, amssymb, amsthm, bm, dcolumn} % Add amsthm here for the proof environment
\usepackage[numbers,sort&compress]{natbib} % Natbib for citations
\usepackage{graphicx} % For high-quality images
\usepackage{multicol} % Multi-column figures/tables
\usepackage{hyperref} % Hyperlinks
\usepackage{listings} % Code listings
\usepackage{authblk} % For structured affiliations
% Define theorem style (optional)
\theoremstyle{plain}
\newtheorem{theorem}{Theorem}
\renewcommand{\qedsymbol}{}

% Custom commands for primes and qubit encoding
\newcommand{\primeQubit}[1]{|p_{#1}\rangle}
\newcommand{\primeGate}[1]{U_{p_{#1}}}

\usepackage{wrapfig}
\usepackage[pscoord]{eso-pic}
\usepackage[fulladjust]{marginnote}
\reversemarginpar

% Typesetting improvements without footnote patching
\usepackage[protrusion=true, expansion=true, tracking=false]{microtype}
\microtypecontext{spacing=nonfrench}

% Line numbers
\usepackage[right]{lineno}

% Text layout - adjust as needed
\raggedright
\setlength{\parindent}{0.5cm}
\textwidth 5.25in 
\textheight 8.75in

% Set double spacing
\usepackage{setspace} 
\doublespacing

% Adjust width for specific content
\usepackage{changepage}

% Adjust caption style
\usepackage[aboveskip=1pt,labelfont=bf,labelsep=period,singlelinecheck=off]{caption}

% Remove brackets from references
\makeatletter
\renewcommand{\@biblabel}[1]{\quad#1.}
\makeatother

% Header, footer, and page numbers
\usepackage{lastpage,fancyhdr}
\pagestyle{fancy}
\fancyhf{}
\fancyhead[L]{Preprint - PrimeAI Enhanced Template}
\fancyfoot[C]{\scriptsize Multiplicity Theory © 2024 Ryan Van Gelder - Citizen Gardens \\ Licensed Under MIT and CC BY-NC-SA 4.0.}
\fancyfoot[R]{Page \thepage\ of \pageref{LastPage}}
\renewcommand{\footrule}{\hrule height 2pt \vspace{2mm}}

\begin{document}

\title{Core Multiplicity Equations}
\author[1]{Ryan Van Gelder}
\author[2]{Nicholas Galioto}
\affil{Citizen Gardens - The Foundation of Multiplicity\\info@citizengardens.org}
\date{\today}

\maketitle
 \begin{multicols}{2}
\begin{singlespace}
\tableofcontents
\end{singlespace}
\end{multicols}
\section{Multiplicity Formula}
\begin{equation}
H \ni \psi \to M(\psi)T(\psi) + f(\psi) = \lambda \psi
\end{equation}
\begin{itemize}
    \item \( H \): Hilbert space.
    \item \( \psi \): State vector.
    \item \( M(\psi) \): Multiplicity operator.
    \item \( T(\psi) \): Coupling tensor.
    \item \( f(\psi) \): Non-linear interaction term.
    \item \( \lambda \): Eigenvalue representing system's response.
\end{itemize}

\section{Time-Dependent Multiplicity Formula}
\begin{equation}
H(t) \ni \psi(t) \to M(t, \psi(t))T(t, \psi(t)) + f(t, \psi(t)) = \lambda(t)\psi(t)
\end{equation}
\begin{itemize}
    \item Time dependencies are added to \( M(t, \psi(t)) \), \( T(t, \psi(t)) \), \( f(t, \psi(t)) \), and \( \lambda(t) \).
    \item Captures dynamic evolution in time-varying systems.
\end{itemize}

\section{General Multiplicity Equation}
\begin{equation}
M(t) = \sum_{i=1}^N \left( \lambda_i \mu_i \cdot e^{i\theta_i(t)} \right) \cdot v_i
\end{equation}
\begin{itemize}
    \item \( \lambda_i \): Eigenvalue.
    \item \( \mu_i \): Multiplicity of the eigenvalue.
    \item \( \theta_i(t) = \omega_i t + \theta_{i0} \): Time-evolving phase.
    \item \( v_i \): Eigenvector.
\end{itemize}

\section{Energy-Based Multiplicity Equation}

The Multiplicity Machine Learning Engine is mathematically represented as:
\begin{equation}
E(t) = (M(t) \cdot S) \otimes T + F(S, H) + \sigma(\omega)
\end{equation}

\begin{itemize}
    \item \textbf{\( E(t) \)}: Energy state of the system over time. It represents the overall dynamics of the system, accounting for all interactions and feedback.
    \item \textbf{\( M(t) \)}: The time-dependent multiplicity operator, capturing the evolving dynamics of the system over time.
    \item \textbf{\( S \)}: The state vector, defining the current state or configuration of the system.
    \item \textbf{\( \otimes \)}: Tensor contraction operation, representing the interaction between the time-dependent multiplicity operator \( M(t) \) and the higher-order coupling tensor \( T \).
    \item \textbf{\( T \)}: The higher-order coupling tensor, capturing complex interactions across different components or variables in the system.
    \item \textbf{\( F(S, H) \)}: A non-linear feedback function that accounts for feedback mechanisms based on the state vector \( S \) and some additional parameters \( H \).
    \item \textbf{\( \sigma(\omega) \)}: A stochastic component introducing randomness or noise into the system, modeled by \( \omega \).
\end{itemize}
\section{Hybrid Multiplicity Equation}
\begin{align}
E(t) &= \sum_{i=1}^N \left[ \left( \lambda_i \mu_i \cdot e^{i\theta_i(t)} \right) \cdot v_i \right] \nonumber \\
&\quad + \left( (M(t, \psi(t)) \cdot S) \otimes T + f(M(t), R(t)) \right) \nonumber \\
&\quad + \lambda(t)\psi(t) + \sigma(\omega)
\end{align}
\begin{itemize}
    \item Tensor interactions:
    \[
    T = \sum_{i,j,k} T_{ijk} \otimes \phi(p_{ijk})
    \]
    \item Stochastic term \( \sigma(\omega) \): Accounts for noise.
\end{itemize}

\section{Dynamic Multiplicity Equation Overview}

The provided Dynamic Multiplicity Equation models the time evolution of the \(k\)-th eigenmode’s intensity \(\rho_k\), incorporating terms for intrinsic dynamics, external inputs, interactions between eigenmodes, and geometric feedback. Below is a detailed breakdown of each term:

\subsection{Equation}
\begin{equation}
\frac{\partial \rho_k}{\partial t} = \alpha_k \rho_k + \beta_k I_k + \gamma_k \sum_j T_{kj} \rho_j + \lambda (\Omega_B + \Omega_{FS})
\end{equation}

\subsection{Terms Explained}

\subsubsection{1. Intrinsic Dynamics: \(\alpha_k \rho_k\)}
\begin{itemize}
    \item \textbf{Description:} This term represents the intrinsic growth or decay of the \(k\)-th eigenmode’s intensity \(\rho_k\).
    \item \textbf{Parameter \(\alpha_k\):} A coefficient controlling the rate of intrinsic dynamics for \(\rho_k\). A positive \(\alpha_k\) suggests exponential growth, while a negative \(\alpha_k\) suggests decay.
    \item \textbf{Role:} Models self-driven processes or inherent tendencies of the eigenmode.
\end{itemize}

\subsubsection{2. External Input: \(\beta_k I_k\)}
\begin{itemize}
    \item \textbf{Description:} Captures the influence of an external input \(I_k\) on the \(k\)-th eigenmode.
    \item \textbf{Parameter \(\beta_k\):} A scaling factor determining the sensitivity of \(\rho_k\) to the input \(I_k\).
    \item \textbf{Role:} Introduces external energy or stimuli, such as environmental changes or external forces.
\end{itemize}

\subsubsection{3. Coupling Between Eigenmodes: \(\gamma_k \sum_j T_{kj} \rho_j\)}
\begin{itemize}
    \item \textbf{Description:} Represents the interaction or coupling between eigenmodes.
    \item \textbf{Parameter \(\gamma_k\):} Determines the strength of the coupling for the \(k\)-th eigenmode.
    \item \textbf{Coupling Matrix \(T_{kj}\):} Encodes the influence of the \(j\)-th eigenmode on the \(k\)-th eigenmode.
    \item \textbf{Role:} Models interconnected dynamics, feedback, and mutual influence across the system.
\end{itemize}

\subsubsection{4. Geometric Feedback: \(\lambda (\Omega_B + \Omega_{FS})\)}
\begin{itemize}
    \item \textbf{Description:} Incorporates feedback from the geometric properties of the system.
    \item \(\Omega_B\): Berry curvature, capturing geometric phase effects and quantum holonomy.
    \item \(\Omega_{FS}\): Fubini-Study metric, related to quantum state fidelity and distances in Hilbert space.
    \item \textbf{Parameter \(\lambda\):} Scales the impact of geometric feedback on the system.
    \item \textbf{Role:} Accounts for non-classical effects and geometric structures that influence the evolution of eigenmodes.
\end{itemize}

\section{The Universal Multiplicity Equation}
The general form of the Universal Multiplicity Equation is expressed as:
\begin{equation}
\frac{\partial \psi_k(t)}{\partial t} = \alpha_k(t)\psi_k + \beta_k(t) \int_0^t I_k(\tau) d\tau + \gamma_k(t) \sum_{j,l} T_{kjl} \psi_j \psi_l + \lambda_k(t) \nabla^2 \psi_k + \eta_k(t) \psi_k^n + \xi_k(t),
\end{equation}
where:
\begin{itemize}
    \item $\psi_k(t)$: The state variable evolving over time.
    \item $\alpha_k(t)$: Growth or decay coefficient, dynamically adjustable.
    \item $\beta_k(t)$: Memory coefficient for historical integration.
    \item $I_k(\tau)$: Input function capturing external influences.
    \item $\gamma_k(t)$: Coupling coefficient for tensor-driven interactions.
    \item $T_{kjl}$: Higher-order tensor encoding multi-scale dependencies.
    \item $\lambda_k(t)$: Quantum potential term for wave-like behavior.
    \item $\eta_k(t)$: Nonlinear self-interaction coefficient.
    \item $n$: Degree of nonlinearity.
    \item $\xi_k(t)$: Stochastic noise term modeling randomness.
\end{itemize}

\section{The Neuromorphic Multiplicity Equation}(NME) has been refined to resolve dimensional inconsistencies and ensure that all terms align with the required dimensions. The corrected form of the NME is expressed as:

\begin{equation}
\frac{\partial \psi_k(t)}{\partial t} = \alpha_k(t) \psi_k + \frac{\beta_k(t)}{T_s} \int_0^t I_k(\tau) \, d\tau + \frac{\gamma_k(t)}{L_s T_s} \sum_{j,l} T_{kjl} \psi_j \psi_l + \frac{\lambda_k(t)}{L_s^2} \nabla^2 \psi_k + \frac{\eta_k(t)}{L_s^{n-1}} \psi_k^n + \xi_k(t),
\end{equation}
where:
\begin{itemize}
    \item \(\psi_k(t)\): State variable with dimensions of \(L\) (length or amplitude).
    \item \(\alpha_k(t)\): Growth or decay coefficient with dimensions \([T^{-1}]\).
    \item \(\beta_k(t)\): Memory coefficient, scaled by the characteristic time \(T_s\), with dimensions \([T^{-1}]\).
    \item \(\int_0^t I_k(\tau) \, d\tau\): Historical integral with dimensions \([L]\).
    \item \(\gamma_k(t)\): Coupling coefficient, scaled by \(L_s T_s\), with dimensions \([L^{-1} T^{-1}]\).
    \item \(T_{kjl}\): Higher-order interaction tensor with dimensions \([1]\).
    \item \(\lambda_k(t)\): Diffusion coefficient, scaled by \(L_s^2\), with dimensions \([L^3 T^{-1}]\).
    \item \(\nabla^2 \psi_k\): Laplacian operator with dimensions \([L^{-1}]\).
    \item \(\eta_k(t)\): Nonlinear interaction coefficient, scaled by \(L_s^{n-1}\), with dimensions \([T^{-1} L^{1-n}]\).
    \item \(\xi_k(t)\): Stochastic noise term with dimensions \([L T^{-1}]\).
    \item \(T_s\): Characteristic time scale, providing temporal scaling.
    \item \(L_s\): Characteristic length scale, providing spatial scaling.
\end{itemize}

This adjustment ensures that each term on the right-hand side (RHS) matches the dimensions of the left-hand side (LHS), which is \([L T^{-1}]\). The inclusion of scaling factors (\(T_s\) and \(L_s\)) harmonizes terms involving time and length dependencies.

\subsection{Interpretation of Adjusted Terms}
\begin{itemize}
    \item \textbf{Feedback and Memory Effects:} The scaling of \(\beta_k(t)\) by \(T_s\) ensures proper weighting of historical integrals to match the dynamics of \(\psi_k(t)\).
    \item \textbf{Tensor-Driven Interactions:} The coupling term is normalized by \(L_s T_s\) to balance interactions involving higher-order tensors.
    \item \textbf{Nonlinear Self-Interactions:} The nonlinear term \(\eta_k(t) \psi_k^n\) is scaled by \(L_s^{n-1}\), allowing emergent behaviors like bifurcations and chaotic dynamics to manifest correctly.
    \item \textbf{Stochastic Noise Handling:} The noise term \(\xi_k(t)\) retains the dimensional consistency needed to model randomness in dynamic systems.
\end{itemize}
This refined equation forms the basis for further theoretical analysis, numerical simulations, and experimental validation across various domains.


\section{Tensor Representation}

To unify various scales and interactions within the Neuromorphic Multiplicity Equation (NME), we represent the system in a tensor-based format, capturing multi-dimensional dependencies and enabling scalability. The tensor-based form of the NME is given by:

\begin{equation}
\frac{\partial \bm{\Psi}(t)}{\partial t} = \bm{A}(t) \bm{\Psi}(t) + \bm{B}(t) \int_0^t \bm{I}(\tau) \, d\tau + \bm{T}(t) \otimes \bm{\Psi}(t) \otimes \bm{\Psi}(t) + \bm{Q}(t) \nabla^2 \bm{\Psi}(t) + \bm{E}(t),
\end{equation}
where:
\begin{itemize}
    \item \(\bm{\Psi}(t)\): State vector representing the system’s evolution over time, with dimensions \([L]\).
    \item \(\bm{A}(t)\): Time-dependent growth or decay tensor, with dimensions \([T^{-1}]\).
    \item \(\bm{B}(t)\): Memory tensor capturing feedback and historical effects, with dimensions \([T^{-1}]\).
    \item \(\bm{I}(\tau)\): Input function tensor capturing external influences, with dimensions \([L T^{-1}]\).
    \item \(\bm{T}(t)\): Higher-order interaction tensor encoding multi-scale dependencies, with dimensions \([L^{-1} T^{-1}]\).
    \item \(\otimes\): Tensor product operator representing multi-dimensional coupling interactions.
    \item \(\bm{Q}(t)\): Quantum-inspired potential tensor for wave-like behaviors, with dimensions \([L^3 T^{-1}]\).
    \item \(\nabla^2 \bm{\Psi}(t)\): Laplacian tensor operator for spatial diffusion, with dimensions \([L^{-1}]\).
    \item \(\bm{E}(t)\): Stochastic noise tensor modeling randomness, with dimensions \([L T^{-1}]\).
\end{itemize}

\subsection{Expanded Terms in Tensor Notation}
The components of the tensor representation are defined as:
\begin{align}
\bm{A}(t) \bm{\Psi}(t) &= \sum_k \alpha_k(t) \Psi_k, \\
\bm{B}(t) \int_0^t \bm{I}(\tau) \, d\tau &= \sum_k \beta_k(t) \int_0^t I_k(\tau) \, d\tau, \\
\bm{T}(t) \otimes \bm{\Psi}(t) \otimes \bm{\Psi}(t) &= \sum_{k,j,l} T_{kjl}(t) \Psi_j \Psi_l, \\
\bm{Q}(t) \nabla^2 \bm{\Psi}(t) &= \sum_k \lambda_k(t) \nabla^2 \Psi_k, \\
\bm{E}(t) &= \sum_k \xi_k(t).
\end{align}

\subsection{Interpretation of Tensor Components}
\begin{itemize}
    \item \textbf{Dynamic Coupling via \(\bm{T}(t)\):} The tensor \(\bm{T}(t)\) encodes complex interactions between different components of the system, enabling multi-dimensional coupling and emergent behaviors.
    \item \textbf{Quantum Potential and Wave Propagation:} The tensor \(\bm{Q}(t)\) governs spatial coherence and wave-like dynamics, leveraging \(\nabla^2\) to model diffusion and tunneling effects.
    \item \textbf{Memory and Feedback Effects:} The integral term \(\bm{B}(t)\) \(\int_0^t \bm{I}(\tau) \, d\tau\) captures historical influences, providing a feedback mechanism for adaptive evolution.
    \item \textbf{Noise Modeling with \(\bm{E}(t)\):} The stochastic tensor \(\bm{E}(t)\) accounts for uncertainties and random fluctuations, ensuring robustness in dynamic environments.
\end{itemize}
\subsection{Dimensional Analysis}
Dimensional consistency ensures all terms align with \([L/T]\), the dimensions of the state vector’s time derivative:
\begin{itemize}
    \item \(\frac{\partial \bm{\Psi}(t)}{\partial t}\): \([L/T]\).
    \item \(\bm{A}(t) \bm{\Psi}(t)\): \(\bm{A}(t) \sim [T^{-1}]\), \(\bm{\Psi}(t) \sim [L]\).
    \item \(\bm{B}(t) \int_0^t \bm{I}(\tau) \, d\tau\): \(\bm{B}(t) \sim [T^{-1}]\), \(\bm{I}(\tau) \sim [L/T]\).
    \item \(\bm{T}(t) \otimes \bm{\Psi}(t) \otimes \bm{\Psi}(t)\): \(\bm{T}(t) \sim [L^{-1} T^{-1}]\), \(\bm{\Psi}(t) \otimes \bm{\Psi}(t) \sim [L^2]\).
    \item \(\bm{Q}(t) \nabla^2 \bm{\Psi}(t)\): \(\bm{Q}(t) \sim [L^3 T^{-1}]\), \(\nabla^2 \bm{\Psi}(t) \sim [L^{-1}]\).
    \item \(\bm{E}(t)\): \([L/T]\).
\end{itemize}
\section{The Multiplicity Constant}
The Multiplicity Constant ($\Lambda_m$) has been a fundamental construct in multiplicity theory, providing a unifying principle for prime-based encoding, tensor networks, and recursive entanglement structures. This work extends the formalization of $\Lambda_m$ by introducing a tensor representation that encodes prime-indexed eigenvalues and quantum entanglement interactions. Additionally, the integration of harmonic oscillatory feedback mechanisms enhances the stability and adaptability of $\Lambda_m$ within recursive quantum systems. These advancements pave the way for applications in quantum-AI systems, neuromorphic intelligence, and topological computation.

The Multiplicity Constant ($\Lambda_m$) serves as a foundational entity in the integration of prime-based encoding with multiplicative tensor networks. Initially formulated as:
\begin{equation}
    \Lambda_m = \lim_{n\to\infty} \sum_{p_i} p_i^{\alpha_i} \left( \xi(p_i) + \psi(p_i, t) \right),
\end{equation}
where $p_i$ are prime-indexed eigenvalues, $\alpha_i$ represents tensor scaling weights, $\xi(p_i)$ accounts for curvature corrections, and $\psi(p_i, t)$ denotes self-referential proof states.

Recent developments have extended $\Lambda_m$ into a tensor-based framework to better accommodate recursive entanglement interactions and dynamic adaptation.
\nocite{*}
\bibliographystyle{unsrt}
\bibliography{references}

\end{document}